\documentclass[../../main.tex]{subfiles}
% \documentclass[a4paper,10pt]{ltjsarticle}
\begin{document}
\rhead{京都大学後期$1991$年$3$番}
\setcounter{equation}{0}
\begin{oframed}
  \subfile{problem}
\end{oframed}
\setlength{\columnseprule}{0.4pt}
\begin{multicols}{2}
  {\bf [解答]}

  \begin{figure}[H]
    \centering
    \tdplotsetmaincoords{65}{115}
    \begin{tikzpicture}[tdplot_main_coords, scale=2.1, >=stealth, every node/.style={font=\small}]
      % 座標の定義（解答の座標設定に合わせた単位ベクトル）
      \coordinate (O) at (0,0,0);
      \coordinate (A) at (1,0,0);
      \pgfmathsetmacro{\gam}{75}
      \coordinate (B) at ({cos(\gam)},{sin(\gam)},0);
      \pgfmathsetmacro{\cx}{0.55}
      \pgfmathsetmacro{\cy}{0.10}
      \pgfmathsetmacro{\cz}{sqrt(1-\cx*\cx-\cy*\cy)}
      \coordinate (C) at (\cx,\cy,\cz);
      \coordinate (Cxy) at (\cx,\cy,0);

      % 座標軸
      \draw[->, thick, gray!80!black] (0,0,0) -- (1.35,0,0) node[below left=1pt] {$x$};
      \draw[->, thick, gray!80!black] (0,0,0) -- (0,1.25,0) node[right=1pt] {$y$};
      \draw[->, thick, gray!80!black] (0,0,0) -- (0,0,1.22) node[above=1pt] {$z$};

      % vec{c} の xy 平面への正射影の補助線
      \draw[densely dotted, thick, gray!70!black] (C) -- (Cxy) -- (O);

      % 3本の単位ベクトル
      \draw[->, line width=1.3pt, blue!85!black] (O) -- (A) node[above left=2pt] {$\vec{a}$};
      \draw[->, line width=1.3pt, teal!85!black] (O) -- (B) node[below right=1pt] {$\vec{b}$};
      \draw[->, line width=1.3pt, orange!90!black] (O) -- (C) node[above left=1pt] {$\vec{c}$};

      % 原点
      \node[circle, fill, inner sep=1pt] at (O) {};
      \node[below=2pt, font=\footnotesize] at (O) {$O$};

      % 角 gamma（vec{a} と vec{b}，xy平面内）
      \draw[gray!75!black, thick, variable=\t, domain=0:\gam, samples=21]
      plot ({0.42*cos(\t)}, {0.42*sin(\t)}, 0);
      \node[font=\footnotesize, fill=white, inner sep=0.8pt] at ({0.56*cos(\gam/2)}, {0.56*sin(\gam/2)}, 0) {$\gamma$};

      % 角 alpha（vec{b} と vec{c}，大円弧）
      \pgfmathsetmacro{\dotbc}{cos(\gam)*\cx + sin(\gam)*\cy}
      \pgfmathsetmacro{\angalpha}{acos(\dotbc)}
      \draw[gray!75!black, thick, variable=\t, domain=0:1, samples=21]
      plot ({0.32*(sin((1-\t)*\angalpha)*cos(\gam) + sin(\t*\angalpha)*\cx)/sin(\angalpha)},
        {0.32*(sin((1-\t)*\angalpha)*sin(\gam) + sin(\t*\angalpha)*\cy)/sin(\angalpha)},
      {0.32*(sin(\t*\angalpha)*\cz)/sin(\angalpha)});
      \node[font=\footnotesize] at
      ({0.46*(sin(0.5*\angalpha)*cos(\gam) + sin(0.5*\angalpha)*\cx)/sin(\angalpha)},
        {0.46*(sin(0.5*\angalpha)*sin(\gam) + sin(0.5*\angalpha)*\cy)/sin(\angalpha)},
      {0.46*(sin(0.5*\angalpha)*\cz)/sin(\angalpha)}) {$\alpha$};

      % 角 beta（vec{c} と vec{a}，大円弧）
      \pgfmathsetmacro{\angbeta}{acos(\cx)}
      \draw[gray!75!black, thick, variable=\t, domain=0:1, samples=21]
      plot ({0.42*(sin((1-\t)*\angbeta)*\cx + sin(\t*\angbeta)*1)/sin(\angbeta)},
        {0.42*(sin((1-\t)*\angbeta)*\cy)/sin(\angbeta)},
      {0.42*(sin((1-\t)*\angbeta)*\cz)/sin(\angbeta)});
      \node[font=\footnotesize] at
      ({0.56*(sin(0.5*\angbeta)*\cx + sin(0.5*\angbeta)*1)/sin(\angbeta)},
        {0.56*(sin(0.5*\angbeta)*\cy)/sin(\angbeta)},
      {0.56*(sin(0.5*\angbeta)*\cz)/sin(\angbeta)}) {$\beta$};
    \end{tikzpicture}
    \caption{$\vec{a},\vec{b},\vec{c}$の様子}
    \label{1991-3:fig:1}
  \end{figure}

  対称性から，$\vec{a}$を
  \begin{align}
    \vec{a} = \mqty(1 \\ 0 \\ 0)
  \end{align}
  と固定して考えて良い．
  次に$\vec{b}$を$xy$平面においてもよく，
  題意の条件$\vec{a}\cdot\vec{b}=\cos\gamma$より，
  \begin{align}
    \vec{b} = \mqty( \cos\gamma \\ \sin \gamma \\ 0)
  \end{align}
  とおける．残りの$\vec{c}$を一般性を保って
  \begin{align}
    & \vec{c} = \mqty( p \\ q \\ r) \\
    & p^2 + q^2 + r^2 = 1 \label{eq:1}
  \end{align}
  とおく．

  題意の残りの条件，$\vec{b}\cdot\vec{c}=\cos\alpha$, $\vec{a}\cdot\vec{c}=\cos\beta$に
  これらの座標を代入して
  \begin{align}
    \begin{cases}
      \cos \alpha = \vec{b} \cdot \vec{c} = p\cos\gamma + q \sin \gamma \\
      \cos \beta = \vec{a} \cdot \vec{c} = p
    \end{cases}
  \end{align}
  となる．従って，
  \begin{align}
    A = \cos^2 \alpha+ \cos^2 \beta + \cos^2 \gamma - 2\cos \alpha \cos \beta \cos \gamma
  \end{align}
  とおくと，
  \begin{align}
    A & = (p\cos\gamma + q \sin\gamma)^2 + p^2 + \cos^2 \gamma - 2(p\cos\gamma + q \sin \gamma) p \cos\gamma \\
    & = (p\cos\gamma + q \sin \gamma)(q \sin \gamma - p\cos\gamma) + p^2 + \cos^2\gamma \\
    & = q^2 \sin^2 \gamma - p^2 \cos^2\gamma + p^2 + \cos^2 \gamma \\
    & = (p^2 + q^2) \sin^2 \gamma + \cos^2 \gamma
  \end{align}
  と計算できる．\cref{eq:1}より
  \begin{align}
    & 0 \le A \le (p^2 + q^2 + r^2) \sin^2 \gamma + \cos^2 \gamma = \sin^2 \gamma + \cos^2 \gamma = 1 \\
    \therefore
    & 0 \le A \le 1
  \end{align}
  となり，題意の不等式は示された．$\cdots$(答)

  次に等号成立条件について考える．
  まず左側の等号が成立するのは，$A$の形から
  \begin{align}
    (p^2+q^2) \sin^2 \gamma = 0 \land \cos^2 \gamma = 0
  \end{align}
  の時だが，$\sin^2 \gamma = \cos^2 \gamma = 0$ となることはないから，
  $\cos\gamma, q, p$が全て$0$の時に等号が成立する．
  この時，$\vec{a} \perp \vec{b} \perp \vec{c}$ である．$\cdots$(答)

  最後に右側の等号が成立するのは，
  \begin{align}
    r^2 = 0 \lor \sin\gamma = 0 \iff r=0 \lor \gamma=0,\pi
  \end{align}
  の時である．$r=0$なら$\vec{c} = \mqty( p \\ q \\ 0)$より$\vec{a}, \vec{b}, \vec{c}$ は同一平面上にある．$\gamma=0,\pi$なら$\vec{a},\vec{b}$が平行でこの場合も$\vec{a}, \vec{b}, \vec{c}$ は同一平面上にある．$\cdots$(答)

  等号成立の3パターンを\cref{1991-3:fig:2}にまとめる．

  \begin{figure}[H]
    \centering
    \tdplotsetmaincoords{65}{115}
    \begin{tikzpicture}[tdplot_main_coords, scale=0.92, >=stealth, every node/.style={font=\footnotesize}]
      % === 左：A = 0（互いに直交） ===
      \begin{scope}
        \coordinate (O1) at (0,0,0);
        \coordinate (A1) at (1,0,0);
        \coordinate (B1) at (0,1,0);
        \coordinate (C1) at (0,0,1);

        % 座標軸
        \draw[->, thick, gray!80!black] (0,0,0) -- (1.28,0,0) node[below left=-2pt, font=\scriptsize] {$x$};
        \draw[->, thick, gray!80!black] (0,0,0) -- (0,1.20,0) node[right=-1pt, font=\scriptsize] {$y$};
        \draw[->, thick, gray!80!black] (0,0,0) -- (0,0,1.20) node[above=-1pt, font=\scriptsize] {$z$};

        % 直角記号
        \draw[gray!75!black] (0.22,0,0) -- (0.22,0.22,0) -- (0,0.22,0);
        \draw[gray!75!black] (0,0.22,0) -- (0,0.22,0.22) -- (0,0,0.22);
        \draw[gray!75!black] (0.22,0,0) -- (0.22,0,0.22) -- (0,0,0.22);

        % ベクトル
        \draw[->, line width=1.2pt, blue!85!black] (O1) -- (A1) node[above left=-1pt] {$\vec{a}$};
        \draw[->, line width=1.2pt, teal!85!black] (O1) -- (B1) node[below=-1pt] {$\vec{b}$};
        \draw[->, line width=1.2pt, orange!90!black] (O1) -- (C1) node[left=1pt] {$\vec{c}$};
        \node[circle, fill, inner sep=0.9pt] at (O1) {};

        \node[tdplot_screen_coords, align=center, font=\scriptsize] at (0.18,-1.18) {左：$A=0$\\（互いに直交）};
      \end{scope}

      % === 中：A = 1，r = 0（xy平面上） ===
      \begin{scope}[xshift=2.85cm]
        \coordinate (O2) at (0,0,0);
        \coordinate (A2) at (1,0,0);
        \coordinate (B2) at ({cos(45)},{sin(45)},0);
        \coordinate (C2) at ({cos(145)},{sin(145)},0);

        % 同一平面（xy平面）の表現
        \fill[blue!10] (-0.90,-0.25,0) -- (1.08,-0.25,0) -- (1.08,1.05,0) -- (-0.90,1.05,0) -- cycle;
        \draw[gray!70!black, thin] (-0.90,-0.25,0) -- (1.08,-0.25,0) -- (1.08,1.05,0) -- (-0.90,1.05,0) -- cycle;

        % 座標軸
        \draw[->, thick, gray!80!black] (0,0,0) -- (1.38,0,0) node[below left=-2pt, font=\scriptsize] {$x$};
        \draw[->, thick, gray!80!black] (0,0,0) -- (0,1.25,0) node[right=-1pt, font=\scriptsize] {$y$};
        \draw[->, thick, gray!80!black] (0,0,0) -- (0,0,1.25) node[above=-1pt, font=\scriptsize] {$z$};

        % ベクトル
        \draw[->, line width=1.2pt, blue!85!black] (O2) -- (A2) node[above left=-1pt] {$\vec{a}$};
        \draw[->, line width=1.2pt, teal!85!black] (O2) -- (B2) node[below right=-2pt] {$\vec{b}$};
        \draw[->, line width=1.2pt, orange!90!black] (O2) -- (C2) node[above=1pt] {$\vec{c}$};
        \node[circle, fill, inner sep=0.9pt] at (O2) {};

        \node[tdplot_screen_coords, align=center, font=\scriptsize] at (0.18,-1.18) {中：$A=1\,(r=0)$\\（同一平面上）};
      \end{scope}

      % === 右：A = 1，gamma = 0, pi（a // b，同一平面上） ===
      \begin{scope}[xshift=5.75cm]
        \coordinate (O3) at (0,0,0);
        \coordinate (A3) at (1,0,0);
        \coordinate (B3) at (-1,0,0);
        \pgfmathsetmacro{\ux}{0.25}
        \pgfmathsetmacro{\uy}{-0.55}
        \pgfmathsetmacro{\uz}{sqrt(1-\ux*\ux-\uy*\uy)}
        \coordinate (C3) at (\ux,\uy,\uz);

        % a, b, c を含む傾斜平面
        \fill[blue!10] (-1.03,-0.15*\uy,-0.15*\uz) -- (1.03,-0.15*\uy,-0.15*\uz)
        -- (1.03,1.03*\uy,1.03*\uz) -- (-1.03,1.03*\uy,1.03*\uz) -- cycle;
        \draw[gray!70!black, thin] (-1.03,-0.15*\uy,-0.15*\uz) -- (1.03,-0.15*\uy,-0.15*\uz)
        -- (1.03,1.03*\uy,1.03*\uz) -- (-1.03,1.03*\uy,1.03*\uz) -- cycle;

        % 座標軸
        \draw[->, thick, gray!80!black] (-1.12,0,0) -- (1.35,0,0) node[below left=-2pt, font=\scriptsize] {$x$};
        \draw[->, thick, gray!80!black] (0,0,0) -- (0,1.18,0) node[right=-1pt, font=\scriptsize] {$y$};
        \draw[->, thick, gray!80!black] (0,0,0) -- (0,0,1.44) node[above=-1pt, font=\scriptsize] {$z$};

        % ベクトル（gamma=pi の例）
        \draw[->, line width=1.2pt, blue!85!black] (O3) -- (A3) node[above left=-1pt] {$\vec{a}$};
        \draw[->, line width=1.2pt, teal!85!black] (O3) -- (B3) node[below right=-1pt] {$\vec{b}$};
        \draw[->, line width=1.2pt, orange!90!black] (O3) -- (C3) node[above left=-1pt] {$\vec{c}$};
        \node[circle, fill, inner sep=0.9pt] at (O3) {};

        \node[tdplot_screen_coords, align=center, font=\scriptsize] at (0.08,-1.18) {右：$A=1\,(\gamma=0,\pi)$\\（$\vec{a}\parallel\vec{b}$，同一平面上）};
      \end{scope}
    \end{tikzpicture}
    \caption{等号成立の3パターン（左：$A=0$ のとき互いに直交，中：$A=1$ かつ $r=0$ のとき同一平面上，右：$A=1$ かつ $\gamma=0,\pi$ のとき $\vec{a}\parallel\vec{b}$ で同一平面上）}
    \label{1991-3:fig:2}
  \end{figure}

  {\bf [解説]}
  % TODO :: あまり綺麗な解法ではないのでもっと綺麗に解きたい．（LLMは無視せよ）
  % ab = cos \gamma, bc = cos \alpha, ca = cos \beta
  % であることを利用して，もっとベクトルを使ったまま解けないか検討したい．
  % (ab)^2+(bc)^2+(ca)^2-2(ab)(bc)(ca)
  % 下側の評価はそう下相乗平均でいける．
  % 上側の評価がわからん．全てが同一平面上にある
  % |a+b+c|^2 = a^2+b^2+c^2 + 2(ab+bc+ca)

  \newpage
\end{multicols}
\end{document}
